\section{治理与未来工作}

\subsection{治理框架}

GR00T 不止研究，也要对齐与监管：

\begin{itemize}
    \item \textbf{Safety board}：每次 release 前必须有人类评估 reward escalation scenarios
    \item \textbf{Usage policy}：定义哪些任务可以 autonomous 执行，哪些必须 human-in-loop
    \item \textbf{Audit trail}：完整记录每个 deployment 的 dataset、hyperparam、model ID
\end{itemize}

\begin{importantbox}{工业级 governance checklist}
1. 是否有激进 failure mode（比如拿起尖锐物）？ 2. 是否提供 override switch？ 3. 是否记录指令 \& response？ 4. 是否有 non-nominal behavior alert？
\end{importantbox}

\subsection{未来研究方向}

Jim Fan 展望三条路径：

\begin{enumerate}
    \item \textbf{LLM + physics solvers}：将 differentiable physics 融入规划
    \item \textbf{Self-supervised manipulation}：用 observation consistency 自动合成 demonstrations
    \item \textbf{Multi-agent robot swarms}：让多个 GR00T 版本协作完成产业流程
\end{enumerate}

\subsection{本章小结}

治理、使用 policy、audit trail 与未来研究一起构成产业落地的安全底座。

